% Part: computability
% Chapter: computability-theory
% Section: equiv-ce-defs

\documentclass[../../../include/open-logic-section]{subfiles}

\begin{document}

\olfileid{cmp}{thy}{eqc}

\olsection[تعريفات المجموعات القابلة للتعداد حسابيًا]{تعريفات متكافئة
  للمجموعات القابلة للتعداد حسابيًا}

تعطي المبرهنة الآتية عددًا من العبارات المهمة المتكافئة لمعنى قابلية
المجموعة للتعداد حسابيًا.

\begin{thm}
\ollabel{thm:ce-equiv}
لتكن $S$ مجموعة من الأعداد الطبيعية. عندئذ تتكافأ العبارات الآتية:
\begin{enumerate}
\item\ollabel{case:ce} $S$ قابلة للتعداد حسابيًا.
\item\ollabel{case:ran-pc} $S$ مدى دالة \emph{جزئية} قابلة للحساب.
\item\ollabel{case:ran-prim} $S$ خالية أو مدى دالة عودية بدائية.
\item\ollabel{case:ce-domain} $S$ \emph{مجال} دالة جزئية قابلة للحساب.
\end{enumerate}
\end{thm}

\begin{explain}
تقول البنود الثلاثة الأولى إنه يمكننا على نحو متكافئ أخذ أي مجموعة غير
خالية قابلة للتعداد حسابيًا معددة بدالة قابلة للحساب، أو بدالة جزئية
قابلة للحساب، أو بدالة عودية بدائية. ويخبرنا البند الرابع أنه إذا كانت
$S$ قابلة للتعداد حسابيًا، فإنه يوجد فهرس~$e$ بحيث
\[
S = \Setabs{x}{\cfind{e}(x) \fdefined}.
\]
وبعبارة أخرى، $S$ هي مجموعة المداخل التي يتوقف عندها حساب
$\cfind{e}$. ولهذا السبب تسمى المجموعات القابلة للتعداد حسابيًا أحيانًا
\emph{شبه قابلة للقرار}: إذا كان عدد في المجموعة حصلت في النهاية على
«نعم»، أما إذا لم يكن فيها فلن تحصل أبدًا على «لا»!{}
\end{explain}

\begin{proof}
لما كانت كل دالة عودية بدائية قابلة للحساب، وكانت كل دالة قابلة للحساب
جزئية قابلة للحساب، فإن \olref{case:ran-prim} تستلزم
\olref{case:ce}، وتستلزم \olref{case:ce} بدورها
~\olref{case:ran-pc}. (ولاحظ أنه إذا كانت $S$ خالية، كانت $S$ مدى الدالة
الجزئية القابلة للحساب التي لا تكون معرفة في أي موضع.) وإذا بينا أن
\olref{case:ran-pc} تستلزم \olref{case:ran-prim}، نكون قد بينا تكافؤ
البنود الثلاثة الأولى.

افترض إذن أن $S$ مدى الدالة الجزئية القابلة للحساب $\cfind{e}$. فإذا
كانت $S$ خالية، انتهينا. وإلا فليكن $a$ أي عنصر من~$S$. وبحسب مبرهنة
كليني للصورة السوية، يمكننا أن نكتب
\[
\cfind{e}(x) = U(\umin{s}{T(e, x, s)}).
\]
وبوجه خاص، تكون $\cfind{e}(x)\fdefined$ و$=y$ إذا وفقط إذا وجد $s$
بحيث تصدق $T(e,x,s)$ ويكون $U(s)=y$. عرف $f(z)$ بـ
\[
f(z) = \begin{cases}
  U((z)_1) & \text{إذا صدقت $T(e, (z)_0, (z)_1)$} \\
  a        & \text{في غير ذلك.}
\end{cases}
\]
عندئذ تكون $f$ عودية بدائية لأن $T$ و$U$ كذلك.
\iftag{TMs}{وبعبارة آلات تورنغ، إذا كان $z$ يرمز زوجًا
$\tuple{(z)_0,(z)_1}$ بحيث تكون $(z)_1$ حسابًا متوقفًا للآلة~$M_e$
عند الدخل $(z)_0$، أعادت $f$ خرج الحساب؛ وإلا أعادت~$a$.}

يلزمنا بيان أن $S$ مدى~$f$؛ أي إنه لكل عدد طبيعي~$y$، يكون $y\in S$
إذا وفقط إذا كان في مدى~$f$. في الاتجاه الأول، افترض أن $y\in S$.
عندئذ يكون $y$ في مدى $\cfind{e}$، ولذلك توجد قيمتان $x$ و~$s$ تصدق
عندهما $T(e,x,s)$ ويكون $U(s)=y$. لكن عندئذ
$y=f(\tuple{x,s})$. وبالعكس، افترض أن $y$ في مدى~$f$. عندئذ إما
$y=a$، وإما أنه يوجد~$z$ بحيث تصدق $T(e,(z)_0,(z)_1)$ ويكون
$U((z)_1)=y$. ولما كانت، في الحالة الأخيرة،
$\cfind{e}((z)_0)\fdefined=y$، فإن $y$ في~$S$ في الحالتين.

(يعني الترميز $\cfind{e}(x)\fdefined=y$ أن «$\cfind{e}(x)$ معرفة
وتساوي $y$». وكان يمكننا أن نستعمل $\cfind{e}(x)=y$ بالقدر نفسه، لكن
السهم الإضافي يفيد أحيانًا في تذكيرنا بأننا نتعامل مع دالة جزئية.)

ولإتمام برهان \olref{thm:ce-equiv}، يكفي أن نبين تكافؤ
\olref{case:ce} و~\olref{case:ce-domain}. لنبين أولًا أن
\olref{case:ce} تستلزم~\olref{case:ce-domain}. افترض أن $S$ مدى دالة
قابلة للحساب~$f$، أي
\[
S = \Setabs{y}{\text{لبعض $x$، } f(x) = y}.
\]
ولتكن
\[
g(y) = \umin{x}{(f(x) = y)}.
\]
عندئذ تكون $g$ دالة جزئية قابلة للحساب، وتكون $g(y)$ معرفة إذا وفقط
إذا كان $f(x)=y$ لبعض~$x$. وبعبارة أخرى، مجال~$g$ هو مدى~$f$.
\iftag{TMs}{وبعبارة آلات تورنغ: إذا أعطيت آلة تورنغ~$F$ تعدد عناصر
~$S$، فلتكن $G$ آلة تورنغ شبه تقرر $S$ بالبحث في مخارج~$F$ لترى هل
عنصر معطى في المجموعة؛ فتتوقف إذا كان فيها، وتواصل البحث إلى الأبد إذا
لم يكن فيها.}{}

وأخيرًا، لبيان أن \olref{case:ce-domain} تستلزم~\olref{case:ce}، افترض
أن $S$ مجال الدالة الجزئية القابلة للحساب~$\cfind{e}$، أي
\[
S = \Setabs{x}{\cfind{e}(x) \fdefined}.
\]
إذا كانت $S$ خالية، انتهينا؛ وإلا فليكن $a$ أي عنصر من~$S$. عرف $f$ بـ
\[
f(z) = \begin{cases}
(z)_0 & \text{إذا صدقت $T(e,(z)_0,(z)_1)$} \\
a & \text{في غير ذلك.}
\end{cases}
\]
وعندئذ، كما سبق، يكون العدد $x$ في مدى~$f$ إذا وفقط إذا كانت
$\cfind{e}(x)\fdefined$، أي إذا وفقط إذا كان $x\in S$.
\iftag{TMs}{وبعبارة آلات تورنغ: إذا أعطيت آلة $M_e$ شبه تقرر $S$،
فعدد عناصر $S$ بالمرور على جميع حسابات آلات تورنغ الممكنة وإعادة
المداخل التي تقابل حسابات متوقفة.}{}
\end{proof}

يوفر البند~\olref{case:ce-domain} من \olref{thm:ce-equiv} طريقة ملائمة
لتعداد المجموعات القابلة للتعداد حسابيًا: لكل~$e$، ليدل $W_e$ على مجال
$\cfind{e}$، أي
\[
W_e = \Setabs{x}{\cfind{e}(x) \fdefined}.
\]
وعندئذ، إذا كانت $A$ أي مجموعة قابلة للتعداد حسابيًا، كان
$A=W_e$ لبعض~$e$.

ويعطي الآتي توصيفًا آخر للمجموعات القابلة للتعداد حسابيًا.

\begin{thm}
\ollabel{thm:exists-char}
تكون المجموعة $S$ قابلة للتعداد حسابيًا إذا وفقط إذا وجدت علاقة قابلة
للحساب $R(x,y)$ بحيث
\[
S = \Setabs{ x }{ \lexists[y][R(x,y)] }.
\]
\end{thm}

\begin{proof}
في الاتجاه الأول، افترض أن $S$ قابلة للتعداد حسابيًا. عندئذ
$S=W_e$ لبعض $e$. ولهذه القيمة من~$e$ يمكننا كتابة $S$ هكذا:
\[
S = \Setabs{ x }{ \lexists[y][T(e, x, y)] }.
\]
وفي الاتجاه العكسي، افترض أن
$S=\Setabs{x}{\lexists[y][R(x,y)]}$. عرف $f$ بـ
\[
f(x) \simeq \umin{y}{R(x, y)}.
\]
عندئذ تكون $f$ جزئية قابلة للحساب، وتكون $S$ مجال~$f$.
\end{proof}


\end{document}
